Befejezésképpen bátran kijelenthetem, hogy a projekt elején leszögezett terveket sikerült implementálni, és fejlesztés közben egyéb hasznos funkcionalitásokat is sikerült hozzáadni a végső eredményhez. Ezáltal a végfelhasználóknak egy-egy funkciódús és jól működő weboldal és applikáció áll a rendelkezésükre, amely nagyban megkönnyíti a tömegközlekedésben való részvételt, bármilyen eszközön is próbálnák elérni ezt. A busztársaságoknak pedig egy egyszerű vezérlőfelület van biztosítva, amely segítségével könnyen és gyorsan tudják kezelni a buszokat, az útvonalakat, a menetrendeket, a jegyárakat és az alkalmazottakat, mindet egy helyen.

Véleményem szerint a Bus4U további idő és munka befektetésével, egy reális piaci rést tud
betölteni, mivel a buszokkal való közlekedés elengedhetetlen a forgalom, és a környezetszennyezés csökkentéséhez, de jelenleg, a rendszer hiányosságai miatt ez egy nem túl népszerű közlekedési eszköz.

\section{További fejlesztési lehetőségek}

Felhasználók szempontjából a következő fontos fejlesztés a buszok telítettségének valós idejű követése lenne, mert ezzel elkerülhető lenne a zsúfoltság a buszokon, és így sokat javulna a tömegközlekedés minősége. Be szeretnénk idővel vezetni a bérletek, illetve különböző kedvezmények kezelésének lehetőségét is, mert ez is növelheti a népszerűséget.

A busztársaságok szempontjából még több lehetőség áll rendelkezésre. Mivel a projektet úgy terveztük, hogy több busztársaság is fel tudja használni, így ez jelen pillanatban általános használatra van optimalizálva. Ennek ellenére egy-egy partner cég kérésére személyre is szabhatnánk a rendszert, egyedi funkciókat és kinézetet bevezetve. Az is hasznos lenne például, ha a cég akár a teljes könyvelését le tudná itt bonyolítani: az alkalmazottak munkaóráit automatikusan vezetni, az autóbuszok tankolásait összesíteni, illetve buszokon lévő POS terminál GPS moduljának segítségével akár a Foaie de parcurs vezetése is automatizálható lenne. A könyveléshez a járművek szervízköltségének számláit is kellene gyűjteni az oldalra, amivel további hasznos kimutatásokat tudnánk generálni.